% automatisch erzeugt von paper/tabellen.py -- nicht von Hand editieren
\begin{table}[htbp]
  \centering
  \footnotesize
  \resizebox{\ifdim\width>\linewidth \linewidth\else \width\fi}{!}{%
  \begin{tabular}{llrrrrrrrrrr}
    \toprule
    Periode & $N$ & \multicolumn{3}{c}{QLIKE} & \multicolumn{3}{c}{MSE ($\times 10^{-7}$)} & \multicolumn{4}{c}{Sharpe} \\
    \cmidrule(lr){3-5}\cmidrule(lr){6-8}\cmidrule(lr){9-12}
     & & Hist. & GARCH & DCC & Hist. & GARCH & DCC & MVP Hist. & MVP GARCH & MVP DCC & 1/N \\
    \midrule
    2003--2007 & 16--16 & -143,5 & -146,8 & -147,3 & 0,28 & 0,23 & 0,23 & 1,190 & 1,146 & 1,158 & 0,845 \\
    2008--2009 & 16--16 & -119,7 & -130,1 & -131,2 & 21,03 & 16,87 & 16,72 & -0,246 & -0,303 & -0,339 & -0,120 \\
    2010--2019 & 16--17 & -146,1 & -147,2 & -147,8 & 0,79 & 0,61 & 0,60 & 0,886 & 0,969 & 0,982 & 0,719 \\
    2020--2021 & 17--17 & -133,1 & -139,5 & -140,3 & 20,18 & 16,67 & 15,88 & 0,436 & 0,391 & 0,417 & 0,828 \\
    2022--2026 & 17--27 & -205,1 & -204,7 & -205,9 & 1,79 & 1,71 & 1,70 & 0,259 & 0,116 & 0,144 & 0,224 \\
    \bottomrule
  \end{tabular}
  }
  \caption{Teilperioden des Basisfalls: Universum, Prognosequalität und Performance}
  \label{tab:strukturbruch_teilperioden}
  \begin{minipage}{\linewidth}\vspace{0.4em}\centering\footnotesize
  $N$ ist die Spannweite der investierbaren Sektoren; QLIKE-Niveaus sind zwischen Perioden mit unterschiedlichem $N$ nicht vergleichbar.
  \end{minipage}
\end{table}
